\section{技术方案}

\subsection{技术选型}

本项目的技术选型围绕三个要求展开。
第一，系统需要能在普通课程项目环境中运行；
第二，模型训练、评估和页面展示要能够复现；
第三，外部服务或大模型权重不可用时，核心分析流程仍应保留。
主要技术栈见表~\ref{tab:technology-selection}。

\begin{table}[H]
  \centering
  \caption{主要技术选型与用途}
  \label{tab:technology-selection}
  \small
  \begin{tabularx}{\textwidth}{C{2.5cm} C{3.4cm} Y}
    \toprule
    层次 & 技术或工具 & 选型理由 \\
    \midrule
    开发语言 & Python 3.10+ &
    便于整合数据处理、模型训练、前端应用和自动化测试 \\
    Web 前端 & Streamlit &
    适合快速构建课程演示型数据应用，支持表格、图表、文件上传和下载 \\
    数据处理 & Pandas、NumPy &
    用于 CSV 清洗、批量统计、训练数据划分和结果导出 \\
    中文处理 & jieba &
    部署轻量，适合课程领域词汇的基础分词和关键词统计 \\
    传统模型 & scikit-learn &
    提供 TF-IDF、Logistic Regression、Naive Bayes、Linear SVM 和评估指标 \\
    深度模型 & PyTorch、Transformers &
    支持 multilingual BERT 微调、本地推理、批量推理和模型保存 \\
    建议生成 & ChatECNU API &
    将结构化分析结果转换为自然语言教学建议；未配置时使用本地模板 \\
    测试工具 & unittest &
    使用 Python 标准库覆盖预处理、模型编排、降级路径和脚本输出 \\
    \bottomrule
  \end{tabularx}
\end{table}

\subsection{系统架构}

系统采用分层结构，将页面交互、NLP 编排、模型推理、规则知识和外部建议服务分开。
前端只调用统一分析接口，不直接依赖某一种模型。
更换情感后端或关闭 LLM 时，页面逻辑不需要同步修改。

\begin{figure}[H]
  \centering
  \resizebox{\textwidth}{!}{%
  \begin{tikzpicture}[
    font=\songti\small,
    layergray/.style={
      draw=ReportLine,
      rounded corners=6pt,
      line width=0.45pt,
      fill=ReportGray
    },
    layerwhite/.style={
      draw=ReportLine,
      rounded corners=6pt,
      line width=0.45pt,
      fill=white
    },
    layerlabel/.style={
      draw=ReportBlue,
      rounded corners=4pt,
      fill=ReportBlue,
      text=white,
      font=\songti\bfseries,
      align=center,
      minimum width=2.45cm,
      minimum height=0.82cm
    },
    module/.style={
      draw=ReportBlue,
      rounded corners=4pt,
      fill=white,
      align=center,
      text width=2.65cm,
      minimum width=3.00cm,
      minimum height=0.92cm,
      inner xsep=4pt,
      inner ysep=3pt,
      line width=0.50pt
    },
    uimodule/.style={
      module,
      fill=ReportLightBlue
    },
    hub/.style={
      draw=ReportBlue,
      rounded corners=5pt,
      fill=ReportLightBlue,
      align=center,
      minimum width=13.40cm,
      minimum height=0.62cm,
      font=\songti\bfseries,
      line width=0.55pt
    },
    resource/.style={
      draw=ReportLine,
      rounded corners=4pt,
      fill=white,
      align=center,
      text width=2.65cm,
      minimum width=3.00cm,
      minimum height=0.82cm,
      inner xsep=4pt,
      inner ysep=3pt,
      line width=0.45pt
    },
    arrow/.style={
      -{Latex[length=2.0mm,width=1.45mm]},
      draw=ReportBlue,
      line width=0.72pt,
      shorten <=2pt,
      shorten >=2pt
    },
    busline/.style={
      draw=ReportBlue,
      line width=0.72pt
    },
    busarrow/.style={
      -{Latex[length=2.0mm,width=1.45mm]},
      draw=ReportBlue,
      line width=0.72pt,
      shorten <=0pt,
      shorten >=0pt
    },
    support/.style={
      -{Latex[length=1.8mm,width=1.25mm]},
      draw=ReportBlue,
      dashed,
      line width=0.52pt,
      shorten <=2pt,
      shorten >=2pt
    }
  ]

    % 三层背景
    \draw[layergray]  (0.00,5.65) rectangle (19.00,7.10);
    \draw[layerwhite] (0.00,1.95) rectangle (19.00,5.35);
    \draw[layergray]  (0.00,-0.05) rectangle (19.00,1.45);

    % 左侧层标签
    \node[layerlabel] at (1.55,6.38) {交互层\\Streamlit};
    \node[layerlabel] at (1.55,3.55) {NLP\\服务层};
    \node[layerlabel] at (1.55,0.70) {模型与\\资源层};

    % 交互层
    \node[uimodule] (single) at (5.20,6.38)
      {单条评价\\即时分析};

    \node[uimodule] (batch) at (9.80,6.38)
      {批量 CSV\\统计分析};

    \node[uimodule] (eval) at (14.40,6.38)
      {固定案例\\模型评估};

    % 统一分析接口
    \node[hub] (hub) at (10.00,4.78)
      {统一分析接口 / NLP Orchestrator};

    % NLP 服务层
    \node[module] (prep) at (4.45,2.88)
      {双语预处理\\情感标准化};

    \node[module] (sent) at (8.60,2.88)
      {情感分析\\BERT + 规则校正};

    \node[module] (evidence) at (12.75,2.88)
      {课程维度识别\\关键词与证据};

    \node[module] (advice) at (16.90,2.88)
      {结构化建议生成\\LLM + 本地模板};

    % 模型与资源层
    \node[resource] (cleanres) at (4.45,0.70)
      {清洗规则\\词表映射};

    \node[resource] (bert) at (8.60,0.70)
      {multilingual BERT\\本地权重};

    \node[resource] (dict) at (12.75,0.70)
      {课程词典\\情感规则};

    \node[resource] (api) at (16.90,0.70)
      {ChatECNU API\\本地兜底};

    % 交互层到统一接口：短竖直箭头
    \draw[arrow] (single.south) -- (single.south |- hub.north);
    \draw[arrow] (batch.south)  -- (batch.south  |- hub.north);
    \draw[arrow] (eval.south)   -- (eval.south   |- hub.north);

    % 统一接口到服务层：总线分发
    \coordinate (busL) at (4.45,3.86);
    \coordinate (busR) at (16.90,3.86);
    \coordinate (busC) at (10.00,3.86);

    \draw[busline] (busL) -- (busR);
    \draw[busline] (hub.south) -- (busC);

    % 分发箭头不使用 shorten，保证与总线、模块边框接上
    \draw[busarrow] (4.45,3.86) -- (prep.north);
    \draw[busarrow] (8.60,3.86) -- (sent.north);
    \draw[busarrow] (12.75,3.86) -- (evidence.north);
    \draw[busarrow] (16.90,3.86) -- (advice.north);

    % 服务层主流程
    \draw[arrow] (prep.east) -- (sent.west);
    \draw[arrow] (sent.east) -- (evidence.west);
    \draw[arrow] (evidence.east) -- (advice.west);

    % 资源层支撑关系
    \draw[support] (cleanres.north) -- (prep.south);
    \draw[support] (bert.north) -- (sent.south);
    \draw[support] (dict.north) -- (evidence.south);
    \draw[support] (api.north) -- (advice.south);

  \end{tikzpicture}}
  \caption{CourseInsight 系统分层架构}
  \label{fig:system-architecture}
\end{figure}

\subsection{双通道数据流}

\begin{figure}[H]
  \centering
  \resizebox{\textwidth}{!}{%
  \begin{tikzpicture}[
    font=\songti\small,
    boxbase/.style={
      rounded corners=4pt,
      line width=0.5pt,
      minimum width=3.40cm,
      text width=3.00cm,
      minimum height=1.15cm,
      align=center,
      inner xsep=4pt,
      inner ysep=4pt
    },
    flowstep/.style={
      boxbase,
      draw=ReportBlue,
      fill=white
    },
    evidencebox/.style={
      boxbase,
      draw=ReportLine,
      fill=ReportGray
    },
    sentimentbox/.style={
      boxbase,
      draw=ReportBlue,
      fill=ReportLightBlue
    },
    arrow/.style={
      -{Latex[length=2.1mm,width=1.5mm]},
      draw=ReportBlue,
      line width=0.8pt,
      shorten <=2pt,
      shorten >=2pt
    },
    branchline/.style={
      draw=ReportBlue,
      line width=0.8pt,
      shorten <=0pt,
      shorten >=0pt
    },
    brancharrow/.style={
      -{Latex[length=2.1mm,width=1.5mm]},
      draw=ReportBlue,
      line width=0.8pt,
      shorten <=0pt,
      shorten >=2pt
    }
  ]

    % 六列等间距节点：列距由 5.10cm 缩短为 4.60cm
    \node[flowstep] (input) at (0.00,0.00)
      {评价输入\\文本 / CSV};

    \node[flowstep] (parse) at (4.60,0.00)
      {文本解析\\原文留存};

    % 上路：原文证据通道
    \node[evidencebox] (raw) at (9.20,1.25)
      {原文表示\\证据保真};

    \node[evidencebox] (extract) at (13.80,1.25)
      {维度定位\\关键词与证据};

    % 下路：标准化情感通道
    \node[sentimentbox] (normalize) at (9.20,-1.25)
      {情感标准化\\缩写与拼写处理};

    \node[sentimentbox] (classify) at (13.80,-1.25)
      {BERT 分类\\规则校正};

    % 汇总与输出
    \node[flowstep] (merge) at (18.40,0.00)
      {结果整合\\情感、维度、证据};

    \node[flowstep] (output) at (23.00,0.00)
      {建议生成\\LLM / 本地模板};

    % 分流与汇流坐标：不显示圆点
    \coordinate (split) at (6.90,0.00);
    \coordinate (join)  at (16.10,0.00);

    % 主流程
    \draw[arrow] (input.east) -- (parse.west);

    % 分流前：直线，无箭头
    \draw[branchline] (parse.east) -- (split);

    % 双通道分流
    \draw[brancharrow] (split) |- (raw.west);
    \draw[brancharrow] (split) |- (normalize.west);

    % 通道内部
    \draw[arrow] (raw.east) -- (extract.west);
    \draw[arrow] (normalize.east) -- (classify.west);

    % 双通道汇流
    \draw[branchline] (extract.east) -| (join);
    \draw[branchline] (classify.east) -| (join);

    % 输出流程
    \draw[brancharrow] (join) -- (merge.west);
    \draw[arrow] (merge.east) -- (output.west);

  \end{tikzpicture}}
  \caption{CourseInsight 双通道数据流}
  \label{fig:data-flow}
\end{figure}

系统对同一条评价保留两种文本表示。
原始文本用于页面展示、课程维度、关键词和证据提取；
标准化文本只用于情感分类。
这样既能提高噪声文本的分类稳定性，也不会因为自动纠错而改变学生原话。

\subsection{数据集与标签设计}

训练数据共 9120 条。
其中 9000 条来自英文 Coursera 课程评论\citep{huggingface_coursera_reviews}；
120 条为人工中文课程评价。
三个情感类别各 3040 条，保持完全均衡。
数据以 Coursera 课程评论为主体；
人工样本用于覆盖中文、中英混合、作业、考试和实验等课堂演示场景。

\begin{table}[H]
  \centering
  \caption{数据来源、用途与结论边界}
  \label{tab:dataset}
  \small
  \begin{tabularx}{\textwidth}{L{3.35cm} C{1.45cm} C{1.85cm} Y C{2.1cm}}
    \toprule
    数据 & 数量 & 语言 & 用途 & 主要量化结论 \\
    \midrule
    Coursera 课程评论 & 9000 & 英文 & 训练、验证和独立测试 & 是 \\
    人工课程评价 & 120 & 中文 & 双语流程验证和课堂演示 & 否 \\
    固定案例 & 12 & 中英混合 & 演示流程验证和回归检查 & 否 \\
    独立压力测试 & 48 & 中英混合 & 复杂表达鲁棒性分析 & 辅助 \\
    \bottomrule
  \end{tabularx}
\end{table}

Coursera 课程评论主要用于英文情感模型训练和量化评估。
人工课程评价样本根据高校课程评价场景构造；
不包含真实学生姓名、学号或其他隐私信息。
报告中的模型性能结论主要依据固定划分的独立测试集；
少量人工样本只用于双语功能验证和测试用例构建。

本项目所说的“中英混合输入支持”，主要指系统输入、预处理、规则识别、课程维度识别和演示流程
能够处理中文、英文及中英混合文本。
模型量化评估主要基于英文课程评论数据。
中文部分目前仍是小规模流程验证，后续需要补充真实中文课程评价数据，才能支撑完整的中文泛化结论。

标签定义为：积极表示主要认可或满意；
消极表示主要抱怨或明确问题；
中性不仅包含态度不明显的评价，也包含建议型评价和正负信息相对平衡的混合评价。
这一定义更符合教学反馈场景。
例如“内容很好，但考试范围不清楚”不应被简单压缩为单一的积极或消极结论。

\subsection{自动降级与运行稳定性}

情感后端默认设置为 \code{AUTO}。
系统优先加载本地 BERT 权重；
若权重缺失、依赖不可用或推理失败，则回退到 TF-IDF Logistic Regression；
如果传统模型文件也不可用，则使用规则模型。
BERT 推理不强制要求 GPU。
当 \code{BERT\_DEVICE=auto} 时，系统会在 CUDA 可用时使用 GPU，否则使用 CPU。
只有强制设置 \code{BERT\_DEVICE=cuda} 且机器没有 CUDA 时，才会触发设备错误。

建议生成同样具有降级路径。
ChatECNU 调用失败或未配置时，本地模板会根据情感、课程维度和证据生成总结、问题与建议。
因此，网络异常只影响外部大模型建议；
不会影响本地情感分类、课程维度识别和关键词提取。